\exercisesection{KRULL DIMENSION OF A RING}

\begin{enumerate}
\def\labelenumi{\arabic{enumi})}
\tightlist
\item
  Let X be a topological space and Y a subspace of X. One says that Y is very dense in X if for every closed subset F of X, F is the closure of F ∩ Y.
\end{enumerate}

\begin{enumerate}
\def\labelenumi{\alph{enumi})}
\item
  For Y to be very dense in X, it is necessary and sufficient that for every nonempty locally closed subset F of X, F ∩ Y be nonempty.
\item
  If Y is very dense in X and Z is a subspace of Y very dense in X, then Z is very dense in X. If Y is very dense in X and F is a locally closed subset of X, then Y ∩ F is very dense in F.
\item
  If Y is very dense in X, one has \(\dim(\mathbf{Y}) = \dim(\mathbf{X})\) .
\end{enumerate}

\begin{enumerate}
\def\labelenumi{\arabic{enumi})}
\setcounter{enumi}{1}
\item
  Let A be a ring and Specmax(A) the set of maximal ideals of A. For Specmax(A) to be very dense in Spec(A), it is necessary and sufficient that A be a Jacobson ring (V, § 3, no. 4). One then has dim Specmax(A) = dim(A).
\item
  Show that if A is a Jacobson ring, A{[}{[}X{]}{]} is not a Jacobson ring. (Compare Specmax(A) and Specmax(A{[}{[}X{]}{]}).)
\item
  \begin{enumerate}
  \def\labelenumii{\alph{enumii})}
  \tightlist
  \item
    Show that one can associate in a unique way to every ordered set E an element of \(\{-\infty\}\cup\mathbf{N}\cup\{+\infty\}\) , denoted dev(E) (deviation of E), such that the following properties are satisfied:
  \end{enumerate}
\end{enumerate}

\(\alpha)\) We have \(\operatorname{dev}(\mathbf{E}) = -\infty\) if and only if the relation \(a \leqslant b\) in \(\mathbf{E}\) implies \(a = b\) .

β) Let n be a positive integer. The deviation of E is ≤ n if and only if, for every infinite decreasing sequence \((a_{k})_{k\in\mathbb{N}}\) of elements of E, one has \(\text{dev}([a_{k}, a_{k-1}]) \leqslant n - 1\) for all sufficiently large k.

\begin{enumerate}
\def\labelenumi{\alph{enumi})}
\setcounter{enumi}{1}
\item
  For one to have \(\text{dev}(\mathbf{E}) \leqslant 0\) , it is necessary and sufficient that every strictly decreasing sequence of elements of E be stationary. One has \(\text{dev}(\mathbf{N}) = 0\) , \(\text{dev}(\mathbf{Z}) = 1\) , \(\text{dev}(\mathbf{Q}) = +\infty\) .
\item
  If E and F are ordered sets and \(f: \mathbf{E} \to \mathbf{F}\) is a strictly increasing map, one has \(\operatorname{dev}(\mathbf{E}) \leqslant \operatorname{dev}(\mathbf{F})\) . We have \(\operatorname{dev}(\mathbf{E} \times \mathbf{F}) = \sup(\operatorname{dev}(\mathbf{E}), \operatorname{dev}(\mathbf{F}))\) .
\item
  Let S(E) be the set of infinite increasing stationary sequences of elements of E, ordered by the product order. Then one has \(\operatorname{dev}(\mathbf{S}(\mathbf{E})) = \operatorname{dev}(\mathbf{E}) + 1\) .
\end{enumerate}

\begin{enumerate}
\def\labelenumi{\arabic{enumi})}
\setcounter{enumi}{4}
\tightlist
\item
  If A is a ring (not necessarily commutative), M a left A-module, one denotes by Kdev(M) the deviation of the set of submodules of M, ordered by inclusion. One sets \(Kdev(A) = Kdev(A_s)\).
\end{enumerate}

\begin{enumerate}
\def\labelenumi{\alph{enumi})}
\tightlist
\item
  If N is a submodule of an A-module M, one has
\end{enumerate}

\[
\operatorname{Kdev} (\mathbf {M}) = \sup (\operatorname{Kdev} (\mathbf {N}), \operatorname{Kdev} (\mathbf {M} / \mathbf {N}))  .
\]

\begin{enumerate}
\def\labelenumi{\alph{enumi})}
\setcounter{enumi}{1}
\item
  Suppose A is commutative. If p and q are two prime ideals of A such that p ⊂ q and p ≠ q, one has dev(A/p) \textgreater{} dev(A/q). (Use the ideals p + Ax, where x is an element of q - p.) Deduce that dim(A) ≤ Kdev(A).
\item
  Suppose A is commutative and noetherian. Prove that one has dim(A) = Kdev(A) (reason by induction on dim(A), assuming that Kdev(B) ≤ dim(B) for every commutative noetherian ring B such that dim(B) \textless{} dim(A); let S be the set of s ∈ A that belong to none of the prime ideals p of A such that dim(A) = dim(A/p); show that, for every finitely generated A-module M such that \(S^{-1}M = 0\) , one has \(\mathrm{Kdev}(\bar{\mathrm{M}}) < \dim(\mathrm{A})\) using the fact that \(S^{-1}A\) is artinian, deduce that \(\mathrm{Kdev}(\mathrm{A}) \leqslant \dim(\mathrm{A})\) . Prove that we have \(\dim(\mathrm{M}) = \mathrm{Kdev}(\mathrm{M})\) for every finitely generated A-module M.
\end{enumerate}

6). a) Let A be a Noetherian (commutative) ring, m an ideal of A contained in the radical of A, and gr(A) the graded ring associated with the m-adic filtration of A. Associate to each ideal of A an ideal of gr(A), and deduce from this that Kdev(A) ≤ Kdev(gr(A)) (cf. Exercise 4, c)). Consequently, one has dim(A) ≤ dim(gr(A)) (Exercise 5, c)).

\begin{enumerate}
\def\labelenumi{\alph{enumi})}
\setcounter{enumi}{1}
\tightlist
\item
  Assume now that m = xA, where x belongs to the radical of A. Show that gr(A) is isomorphic to a quotient of the graded algebra (A/m)[T]. Show that the ordered set F of graded ideals of (A/m) {[}T{]} is isomorphic to the set of increasing sequences of ideals of A/m; deduce from this (exerc. 4, d)) that one has dev(F) = Kdev(A/(x)) + 1, then dim(A) ≤ dim(A/m) + 1.
\end{enumerate}

\begin{enumerate}
\def\labelenumi{\arabic{enumi})}
\setcounter{enumi}{6}
\tightlist
\item
  Let A be the ring of germs at 0 of functions of class \(C^\infty\) on \(\mathbf{R}_{+}\). We propose to show that the ring A is of infinite dimension. We denote by C the algebra of functions of class \(C^\infty\) on \(\mathbf{R}_{+}\) , and let F be the ideal formed by the functions vanishing in a neighborhood of 0, so that \(A = C / F\) .
\end{enumerate}

\begin{enumerate}
\def\labelenumi{\alph{enumi})}
\item
  Let \((x_{n})_{n\in \mathbb{N}}\) be a strictly decreasing sequence, tending to 0, of elements of \(\mathbf{R}_{+}\) , and let \(\mathfrak{U}\) be an ultrafilter on \(\mathbf{N}\) finer than the filter of complements of finite sets. Show that the set \(\mathbf{J}_1\) of \(f\in \mathbf{C}\) such that the set of \(n\in \mathbf{N}\) such that \(f(x_{n}) = 0\) belongs to \(\mathfrak{U}\) is a prime ideal of \(\mathbf{C}\) .
\item
  We denote by \(\tau_{n}(f)\) the function \(x\mapsto f(x + x_{n})\) . Show that if J is a prime ideal of C, the set \(\varphi(J)\) of \(f\) such that the set of \(n\in\mathbf{N}\) with \(\tau_{n}f\in\mathbf{J}\) belongs to \(\mathfrak{U}\) is a prime ideal of C.
\item
  One defines a sequence \(\mathbf{I}_n\) of ideals of C by \(\mathbf{I}_0 = \mathbf{J}_1\) , \(\mathbf{I}_{n+1} = \varphi(\mathbf{I}_n)\) . Show that \((\mathbf{I}_n)_{n\in\mathbb{N}}\) is a strictly decreasing sequence of prime ideals of C containing F.
\end{enumerate}

\begin{enumerate}
\def\labelenumi{\arabic{enumi})}
\setcounter{enumi}{7}
\item
  Let X be a topological space. Show that the map \(x \mapsto \dim_x(X)\) from X to \(\overline{\mathbf{R}}\) is upper semicontinuous.
\item
  Let A be a ring, p a prime ideal of A, and M a finitely generated A-module. Show that one has:
\end{enumerate}

\[
\dim_ {\mathbf {A} _ {\mathfrak {p}}} (\mathbf {M} _ {\mathfrak {p}}) + \dim_ {\mathbf {A} / \mathfrak {p}} (\mathbf {M} / \mathfrak {p} \mathbf {M}) \leqslant \dim_ {\mathbf {A}} (\mathbf {M}).
\]

\begin{enumerate}
\def\labelenumi{\arabic{enumi})}
\setcounter{enumi}{9}
\tightlist
\item
  A ring A is said to be equidimensional if all irreducible components of Spec(A) have the same dimension. Let A be a ring and n an integer. The following conditions are equivalent:
\end{enumerate}

\begin{enumerate}
\def\labelenumi{\alph{enumi})}
\tightlist
\item
  every chain of prime ideals of A is contained in a chain of length n;
\item
  for every maximal ideal m of A, the ring A \(_{m}\) is equidimensional, catenary, and of dimension n.
\end{enumerate}

\begin{enumerate}
\def\labelenumi{\arabic{enumi})}
\setcounter{enumi}{10}
\tightlist
\item
  Let \((\mathbf{A}_{i}, f_{ij})\) be an inductive system of rings relative to a filtered ordered set I, with inductive limit A. Show that one has \(\dim(\mathbf{A}) \leqslant \lim.\inf\dim(\mathbf{A}_{i})\) and that one has equality when the homomorphisms \(f_{ij}\) are faithfully flat. Give an example where the homomorphisms \(f_{ij}\) are flat and where one has \(\dim(\mathbf{A}) < \lim.\inf\dim(\mathbf{A}_{i})\) (take I = N, \(A_{n} = Z[n^{-1}]\) , A = Q).
\end{enumerate}

Let A be a ring, and let I and J be two ideals of A. Suppose that for every maximal ideal m of A, we have IA \(_{m}\) = JA \(_{m}\) . Deduce from this that I = J.

\begin{enumerate}
\def\labelenumi{\alph{enumi})}
\setcounter{enumi}{1}
\item
  Let A be a ring such that, for every maximal ideal m of A, the ring \(\mathrm{A}_{\mathfrak{m}}\) is Noetherian and such that, for every nonzero element \(f\) of A, there exist only finitely many maximal ideals of A containing \(f\) . Then A is Noetherian. (Let \(a\) be an ideal of A. We will first show that there exists a finite family \((x_1, ..., x_n)\) of elements of \(a\) such that the maximal ideals of A that contain \(\{x_1, ..., x_n\}\) contain \(a\). We will then construct a finite family of elements \(y_1, ..., y_l\) such that for every maximal ideal m of A, \(\{y_1, ..., y_l\}\) generates the ideal \(a\mathrm{A}_{\mathfrak{m}}\) of \(\mathrm{A}_{\mathfrak{m}}\) . We conclude by using a).
\item
  Let A be a semi-local ring such that, for every maximal ideal m, the ring \(A_{m}\) is Noetherian. Then A is Noetherian.
\end{enumerate}

\begin{enumerate}
\def\labelenumi{\arabic{enumi})}
\setcounter{enumi}{12}
\item
  Let K be a field, \((n_{r})_{r\geqslant1}\) a strictly increasing sequence of integers \textgreater0; for every integer \(i\geqslant0\) , denote by \(p_{i}\) the ideal of the ring \(\mathrm{K}[(\mathrm{X}_{j})_{j\in\mathbb{N}}]\) generated by the \(X_{j}\) for \(n_{1}+\cdots+n_{i}\leqslant j<n_{1}+\cdots+n_{i}+n_{i+1}\) , and let S be the set of elements of \(\mathrm{K}[(\mathrm{X}_{j})]\) that belong to none of the \(p_{i}\). Then the ring \(S^{-1}K[(X_{j})]\) is Noetherian and of infinite dimension (use Exercise 12 to show that \(S^{-1}K[(X_{j})]\) is Noetherian) \(^{1}\) .
\item
  Let V be a discrete valuation ring, and let \(\pi\) be a uniformizer of V. In the ring V{[}T{]}, the ideal \(m_1 = (\pi T - 1)\) is maximal and of height 1, the ideal \(m_2 = (\pi) + (T)\) is maximal and of height 2. The fields V{[}T{]}/ \(m_1\) and V{[}T{]}/ \(m_2\) are isomorphic to the field of fractions of V and to the residue field of V, respectively. One has \(\dim(\mathrm{V}[T]) = 2\), but \(\dim(\mathrm{V}[T]/m_1) + \dim(\mathrm{V}[T]_{m_1}) = 1\) and \(\dim(\mathrm{gr}_{m_1}(\mathrm{V}[T])) = 1\).
\end{enumerate}

With the notation of the preceding exercise, suppose that the residue field and the field of fractions of the ring V are isomorphic (this is the case, for example, when \(\mathrm{V} = k[[\mathrm{U}]]\) , the field \(k\) being the field of fractions of a ring of formal power series in infinitely many indeterminates with coefficients in a field). Let \(\sigma\) be an isomorphism of \(\mathrm{V[T] / m_1}\) onto \(\mathrm{V[T] / m_2}\) . Let C denote the subring of \(\mathrm{V[T]} = \mathrm{E}\) consisting of the elements of E whose classes modulo \(m_1\) and \(m_2\) correspond under \(\sigma\) . Show that C is Noetherian and that \(m_1m_2 = m_1 \cap m_2\) is a maximal ideal of C. Show that E is the integral closure of C (noting that \(\mathrm{E} = \mathrm{C} + \mathrm{C}\cdot(\pi T - 1)\) and that \((\pi T - 1)^2 + (\pi T - 1)\) belongs to C).

Let \(\mathring{\mathrm{A}}\) be the local ring \(\mathbf{C}_{\mathfrak{m}_1\mathfrak{m}_2}\) . Then A is an integral domain of dimension 2, hence catenary; the integral closure \(\mathbf{B} = \mathbf{E} \otimes_{\mathbb{C}} \mathbf{A}\) of \(\mathring{\mathrm{A}}\) is a semi-local Noetherian ring having exactly two maximal ideals \(\mathfrak{n}_1\) and \(\mathfrak{n}_2\) , and we have \(\dim(\mathbf{B}_{\mathfrak{n}_1}) = 1\) , \(\dim(\mathbf{B}_{\mathfrak{n}_2}) = 2\) .

¶ 16) We keep the notation of the preceding exercise, and we identify B with the quotient of the polynomial ring A{[}U{]} by the prime ideal \(\mathfrak{p}\) generated by \(U^2 + U - \pi T(\pi T - 1)\). Let \(\mathfrak{q}_i\) be the ideal of A{[}U{]} such that \(\mathfrak{n}_i = \mathfrak{q}_i/\mathfrak{p}\). Then \(\operatorname{ht}(\mathfrak{q}_i) = 3\), \(\operatorname{ht}(\mathfrak{p}) = 1\), and the chain \(\mathfrak{p} \subset \mathfrak{q}_1\) is saturated. In particular, A{[}U{]} and A{[}U{]}\(_{\mathfrak{q}_1}\) are not catenary.

\begin{enumerate}
\def\labelenumi{\arabic{enumi})}
\setcounter{enumi}{16}
\tightlist
\item
  \begin{enumerate}
  \def\labelenumii{\alph{enumii})}
  \tightlist
  \item
    Let A be an integral domain and \(f \in \mathrm{A}[\mathrm{T}], f \neq 0\). Show that there exists a maximal ideal m of \(\mathrm{A}[\mathrm{T}]\) such that \(f \notin \mathfrak{m}\) (take a maximal ideal containing T.f-1).
  \end{enumerate}
\end{enumerate}

\begin{enumerate}
\def\labelenumi{\alph{enumi})}
\setcounter{enumi}{1}
\tightlist
\item
  Let A be a ring. Show that one has dim(Specmax(A{[}T{]})) \(\geqslant\) dim(A) + 1. (Let \(\mathfrak{p}_0\subset \mathfrak{p}_1\subset \cdots \subset \mathfrak{p}_n\) be a chain of prime ideals of A. Consider the chain
\end{enumerate}

\[
\mathfrak {p} _ {0} [ \mathrm{T} ] \subset \mathfrak {p} _ {1} [ \mathrm{T} ] \subset \mathfrak {p} _ {2} [ \mathrm{T} ] \dots \subset \mathfrak {p} _ {n} [ \mathrm{T} ] \subset (\mathfrak {p} _ {n}, \mathrm{T})
\]

of prime ideals of A{[}T{]} and let us set

\[
\left\{ \begin{array}{l} \mathrm{U} _ {i} = \mathrm{V} (\mathfrak {p} _ {i} [ \mathrm{T} ]) \cap \operatorname{Specmax} (\mathrm{A} [ \mathrm{T} ]), \quad 0 \leqslant i \leqslant n \\ \mathrm{U} _ {n + 1} = \mathrm{V} (\mathfrak {p} _ {n}, \mathrm{T}) \cap \operatorname{Specmax} (\mathrm{A} [ \mathrm{T} ]). \end{array} \right.
\]

Show, using \(a\) ) that \((\mathrm{U}_i)_{0 \leqslant i \leqslant n+1}\) is a strictly increasing sequence of irreducible closed subsets of Specmax (A{[}T{]}).

\begin{enumerate}
\def\labelenumi{\alph{enumi})}
\setcounter{enumi}{2}
\tightlist
\item
  Let \(k\) be a field, and let S be the subset of \(k[\mathrm{X}, \mathrm{Y}]\) formed by the elements of the form \(1 + \mathrm{X}g(\mathrm{X}, \mathrm{Y})\) , and let \(\mathrm{A} = \mathrm{S}^{-1}k[\mathrm{X}, \mathrm{Y}]\) .
\end{enumerate}

Show that one has dim(Specmax(A)) = 1 and dim(Spec(A)) = 2. Generalize to the case of several indeterminates.

\begin{enumerate}
\def\labelenumi{\alph{enumi})}
\setcounter{enumi}{3}
\item
  Show that for every pair of integers \(0 \leqslant n < m\) , there exists a Noetherian ring A such that \(\dim(\text{Specmax}(A)) = n\) and \(\dim(\text{Specmax}(A[T])) = m\) .
\item
  Let A be the Noetherian ring defined in exerc. 13. Show that one has dim(Specmax(A)) = 1 and dim(Specmax(A{[}T{]})) = \(\infty\) .
\end{enumerate}

